\chapter{The Semiconductor: Fabs, Logic, and DRAM}
\chaptermark{The Semiconductor}
\label{ch:semi}

The third element is the physical base. Models can be open-sourced and architectures standardized overnight; fabs cannot. This is where the West's chokepoint strategy concentrated---EUV denial since 2019, successive equipment and HBM controls, entity lists---and therefore where the empirical test of \emph{chokepoint decay} (Chapter~\ref{ch:playbook}) is sharpest. The evidence through mid-2026: the chokepoints are decaying faster than they are being replaced.

\section{Logic: 5\,nm without EUV}

SMIC crossed the line that export controls were designed to hold twice. The Huawei Mate~60's 7\,nm Kirin 9000S (August 2023) was the first shock; by December 2025, TechInsights teardowns confirmed volume production on SMIC's N+3 node---a scaled evolution of its 7\,nm-class technology that TechInsights describes as a key indicator of how close SMIC is to a true 5\,nm-equivalent node without EUV, while remaining ``significantly less scaled than leading commercial 5\,nm'' from TSMC or Samsung [V]~\cite{smic-5nm,techinsights-n3}. Capability is demonstrated; TSMC-equivalent economics are not, and the comparison metrics that would settle it---density, yield by die size, wafer cost, cycle time---remain undisclosed [A]. The costs are real: multi-patterning burns wafer capacity, N+2 yields are estimated at 40--55\% for large dies (Cambricon's biggest reportedly ${\sim}20\%$), and margins show the strain~\cite{yields,cambricon}. But the strategic arithmetic is not yield percentage; it is wafers times yield times time. SMIC's advanced-node capacity is climbing from ${\sim}45{,}000$ wafers/month (2025) toward 60,000 (2026) and 80,000 (2027)~\cite{semianalysis-ascend}; SMIC posted record 2025 revenue of \$9.3B at 93.5\% utilization~\cite{smic-2025}. Around it, Huawei's ``shadow fab'' network---PXW, SwaySure, Pengjin, SiEn, and others identified by open-source analysis, roughly eleven fabs of which eight are 300\,mm---adds capacity outside the visible SMIC envelope~\cite{shadow-fabs}. Huawei's die bank (over 2.9 million Ascend dies procured from TSMC before enforcement closed the loophole, plus $\sim$13 million HBM stacks stockpiled) bridges the gap while domestic capacity ramps~\cite{semianalysis-ascend}.

Below the leading edge, the mature-node story is already solar-shaped: China holds 37\% of global 22--40\,nm output in 2026 (42\% projected by 2028) and drives nearly 70\% of global 300\,mm mature-node capacity expansion, with trailing-edge price pressure forcing Taiwanese incumbents through the familiar sequence of supplier squeezes and order migration~\cite{mature-node}. Every power stage, microcontroller, and sensor of the EV and robotics stack rides on these nodes---and recall (Chapter~\ref{ch:ev}) that Chinese automakers are directed toward 100\% domestic chips by 2027~\cite{auto-chips}.

\section{DRAM: CXMT and the swing-supplier position}

Memory is the quiet revolution. CXMT grew from ${\sim}70{,}000$ wafers/month in 2022 to 200,000--300,000 by mid-2026, targeting 600,000 via new Hefei and Shanghai fabs; independent modeling puts it within ${\sim}25{,}000$ wafers/month of Micron by end-2026---which would make China's national champion the world's fourth, and trajectory-wise third, DRAM producer~\cite{cxmt-capacity,citrini}. Its DDR5 yields reached 80\%+; it publicly launched DDR5-8000 and LPDDR5X-10667 in November 2025; and its July 2026 Shanghai IPO raised \$8.6B and closed +466\% at a ${\sim}\$460$B market capitalization---instantly the largest A-share company~\cite{cxmt-ddr5,cxmt-ipo,cxmt-reuters}. Three disciplines keep that number honest. First, only 6.73\% of shares floated, so the print measures state-channelled investor appetite in a thin market rather than industrial value---first-day market capitalization is a \emph{sentiment} indicator and nothing more [V]~\cite{cxmt-reuters}. The industrial indicators are the ones to track instead: IPO proceeds and their allocation, wafer capacity by node, revenue, gross margin, yield, customer qualification, and the capex plan. Second, the prospectus disclosed \emph{no dedicated commercial HBM investment}, which sits awkwardly beside the narrative---this book's included---that CXMT HBM3 is an imminent strategic bridge. The contradiction should be stated rather than smoothed: either HBM development is funded outside the disclosed lines, buried inside general R\&D, or less mature than the reporting implies. All three readings are live, and the resolution is a disclosure this book does not have [A]. Third, the swing-supplier evidence (CXMT server modules pricing above Samsung's) is a moment in a panic market, not proof of durable pricing power: qualification status, channel inventory, specification, shortage timing, and regional procurement all move that number.

That last fact is the essential one, and it deserves the full accounting, because domestic HBM is now the single binding physical constraint on the compute element---and the most important open variable in this book.

\subsection{The HBM ledger}
\label{sec:hbm-ledger}

The state of Chinese HBM as of mid-2026, stated carefully by evidence grade. \emph{In production}: CXMT HBM2 since 2H24, 8-high HBM2E since 1H25---roughly three to four years behind SK~hynix, on a $\sim$16\,nm-class cell with a $>$30\% cost-per-bit disadvantage, and today under 2\% of CXMT's $\sim$265,000 monthly wafer starts~\cite{chinatalk-hbm,cxmt-ipo}. \emph{Sampled}: CXMT delivered HBM3 samples to Huawei in October 2025, with large-scale HBM3 manufacturing slated for 2026 and HBM3E for 2027; Korean-sourced reports have CXMT dedicating $\sim$20\% of capacity to an HBM3 line~\cite{cxmt-hbm3}. \emph{Announced}: Huawei's Ascend 950PR ships in 2026 with self-developed ``HiBL~1.0'' HBM (128\,GB, 1.6\,TB/s) and the 950DT with ``HiZQ~2.0'' (144\,GB, 4\,TB/s)---whose DRAM dies are fabbed by an undisclosed partner, with CXMT and the entity-listed Wuhan Xinxin (XMC, $\sim$3,000 HBM wafers/month with SJSemi packaging) the documented candidates~\cite{huawei-roadmap,xmc-hbm}. \emph{Claimed}: Chinese state media assert LineShine runs on ``China's first domestically developed HBM''---which would make the world's No.~1 supercomputer the first at-scale deployment of Chinese HBM---while independent analysts note the supplier is undisclosed and speculate CXMT HBM2E~\cite{globaltimes-lineshine,lockwood}. \emph{Not yet verified}: no public teardown has confirmed Chinese-fabbed DRAM dies in a commercially shipping accelerator; teardowns of shipping Ascend 910Cs still find Samsung and SK~hynix stacks~\cite{hbm-teardown}.

Those foreign stacks are the depleting bridge: Huawei accumulated $\sim$13 million HBM stacks (including $\sim$11.4M from Samsung, $\sim$7M of them shipped in the December 2024 window between the BIS rule's announcement and its enforcement), and by mid-2026 that stockpile is assessed as effectively exhausted---making HBM, not SMIC logic, the binding constraint on Ascend output ($\sim$2M domestic stacks projected for 2026, sufficient for only $\sim$250--300k Ascend 910C-class packages; more bullish estimates reach $\sim$7M)~\cite{semianalysis-ascend,substrate}. The advanced-packaging leg, by contrast, is scaling fast: Innotron's \$2.4B TSV/stacking facility, Tongfu's RMB 4.4B raise, over RMB 27B poured into packaging expansion in H1 2026, and a domestic HBM assembly toolchain (Naura TSV etch, Maxwell hybrid bonding) targeted at end-2026 HBM3~\cite{packaging,hbm-toolchain}. Washington noticed the champion late: CXMT reached the DoD's 1260H list only in June 2026 and remains off the BIS Entity List---after the IPO, after the LPDDR6 verification, after the HBM3 samples~\cite{cxmt-ipo}.

The strategic read is double-tracked, and Chapter~\ref{ch:compute}'s template (Section~\ref{sec:template}) is its exact mirror. On the HBM track, China is three-plus years behind and sprinting---independent assessments (TrendForce) still place domestic HBM in early verification with limited near-term global impact [A]~\cite{trendforce-hbm}---with the inference-centric machines (Ascend 950, LX2-class sockets) as the guaranteed captive demand. On the commodity track, CXMT is at the \emph{world's leading edge} of LPDDR6---sampling concurrent with SK~hynix, mass production targeted 2H26 [R]---feeding the training-centric, capacity-first architecture designed to need far less HBM~\cite{cxmt-lpddr6}. The precise formulation matters, because the looser version of it is false: China has not made HBM irrelevant; it is engineering a position in which HBM ceases to be the \emph{sole} architectural path to frontier-scale AI, while racing to supply it domestically anyway. Section~\ref{sec:hbm4-spread} then adds a complication the ledger above cannot absorb without revision: during the generation China is chasing, the thing being chased stopped being a single specification.

\subsection{The 2026 memory wall is priced, not clocked}
\label{sec:memory-wall}

Meanwhile the 2025--2026 global memory shortage---DRAM contract prices up 50--60\% in quarters, driven by AI demand and HBM's $3\times$ wafer intensity---handed CXMT the swing-supplier role: its 64\,GB server DDR5 modules now price \emph{above} Samsung's~\cite{dram-shortage,cxmt-pricing}. Within twelve months, the narrative flipped from ``Chinese dumping'' to Chinese pricing power, and that evidence remains a moment in a panic market rather than proof of durable pricing power [A]. YMTC's 294-layer NAND and the revival of Entity-Listed Fujian Jinhua complete the memory picture~\cite{ymtc}. By mid-2026, though, ``shortage'' had become the wrong word for what the series show, and the correction matters here because it changes which side of this book's argument the memory market sits on.

The most complete public accounting is Jim Handy's tutorial at Hot Chips 2026, and it should be read with its conflict on the table---the deck itself discloses that the data underpin a paid report, so these are an analyst's series and not an audited one [A]~\cite{hc2026-memory-wall}. What they show is a price event, not a capacity event. The lowest spot DRAM price is up $7.1\times$ and NAND $6.9\times$ from the mid-2025 trough. Combined DRAM and NAND revenue at the five major suppliers has gone from a ${\sim}\$15$B quarterly trough in 2Q23 to ${\sim}\$180$B in 2Q26---a twelvefold move that is almost entirely price and almost not at all bits. Hyperscale property and equipment spending rose over the same era from ${\sim}\$30$B a quarter in 1Q21 to ${\sim}\$170$B a quarter in 2Q26. And the supply side did not move: DRAM wafer starts have been flat for more than a decade, and new capacity takes over two years to install, so no supplier can answer this on the clock even if every one of them has decided to. HBM is the mechanism---roughly $3\times$ the die area per bit of DDR, so each stack sold withdraws several stacks' worth of commodity bits from the market---and the consequence surfaces not as an allocation queue but as a price. Handy's summary slide is a photograph of banknotes.

The two price series in the preceding paragraphs measure different things, and both are needed. Contract prices up 50--60\% in a quarter is what a qualified customer with a long-term agreement pays; a $7.1\times$ move off the trough in the lowest spot print is what a buyer without one pays---and the buyers without one are precisely the marginal actors this book cares about: systems integrators, second-tier clouds, national laboratories, and anyone standing up capacity that was not in a 2024 forecast. Two vendor figures fix the magnitude at the system level. Samsung's own accounting puts memory at 63\% of AI-chip component spend in 4Q25, up from 52\% in 1Q24 [P]. Micron's is more arresting: an eight-stack HBM4 package carries 12{,}344\,mm$^2$ of memory silicon, more than $8\times$ the die area of a typical GPU [P]~\cite{hc2026-hbm4-spread}. On an AI accelerator the memory is now the majority of the silicon and the majority of the bill of materials; memory price sets system price, and the logic vendor's margin is what absorbs the residual.

This cuts both ways for the argument of this book, and the inconvenient direction should be stated before the convenient one. It raises the cost of the West's HBM-centric architecture, which is the capacity-first thesis of Chapter~\ref{ch:compute} being priced by the market rather than argued on a slide: if memory is 63\% of the bill and the HBM fraction of it is the most expensive memory ever manufactured, then an architecture substituting capacity for bandwidth is buying the cheaper end of a market whose expensive end is running away. But it raises the cost of the cheaper end too. The ${\sim}$2.3\,TB commodity socket of Section~\ref{sec:lpddr6} is denominated in exactly the bits whose spot price rose $7.1\times$, and the appliance of Section~\ref{sec:appliance-cost} in the NAND that rose $6.9\times$. \emph{A memory price shock is not unambiguously good news for the side that chose to buy more memory of a cheaper kind}---it is good news only insofar as the cheap kind stays cheaper in \emph{relative} terms, which the HBM die-area arithmetic supports and the flat wafer starts do not guarantee. Where this book prices a capacity-first machine it does so on a pre-surge basis, and that basis is now a modeling assumption rather than an observation [M].

What the shock does do for China is narrower than a price move and more durable. First, it converts CXMT's capacity into revenue at the moment that revenue funds the next tranche of capacity; the IPO proceeds and the price cycle arriving in the same quarter is the luckiest piece of timing in the domestic memory program. Second, the two-year installation lag binds everyone equally, and CXMT's wafer additions are among the largest single tranches of new commodity DRAM capacity being installed anywhere---in a market that cannot be answered on the clock, the fastest-growing marginal supplier gains share by default. Third, and this is the term the controls do not reach: the shortage is a consequence of the West's own architectural choice. HBM's die-area intensity is what withdrew the commodity bits, and the country that elected not to build HBM at scale is the country whose fab capacity is not being consumed by it.

\section{HBM4 stopped being a specification}
\label{sec:hbm4-spread}

The ledger of Section~\ref{sec:hbm-ledger} measures a gap in years. That measurement presumes the thing being chased holds still enough to be chased, and the 2026 disclosure record---the three incumbent suppliers' own decks---shows that within one generation it stopped doing so. Two changes arrived together, and between them they alter what ``three to four years behind'' means. The conference record in full is Appendix~\ref{app:hotchips}; what follows is the part bearing on this chapter.

\subsection{A $1.5\times$ spread inside one nominal generation}

The same JEDEC HBM4 generation is quoted by SK~hynix at 8\,Gb/s and 2.0\,TB/s per stack, by Micron at 11\,Gb/s, and by Samsung at 3.0\,TB/s---a $1.5\times$ spread in per-stack bandwidth within one nominal generation [P]~\cite{hc2026-hbm4-spread}. A bill of materials specifying ``HBM4'' therefore does not specify a part, and the generation's two flagship accelerators made opposite choices from it. AMD's MI455X bought the slow bin and compensated with stack count: twelve stacks, 432\,GB, 23.3\,TB/s, an implied ${\sim}7.6$\,GT/s. NVIDIA's Rubin bought the fast bin and stayed at eight: 288\,GB, 22.2\,TB/s, an implied ${\sim}10.8$\,GT/s---having re-specified onto the faster bin partway through the program, with the per-GPU envelope rising from 1.8 to 2.3\,kW. The re-specification is legible in NVIDIA's own record, since the March 2026 launch material described the same 288\,GB part at ${\sim}13$\,TB/s [P]~\cite{nvidia-rubin-pod,vr200-power}. The two designs land within 5\% of each other on aggregate bandwidth by entirely different routes---one buying capacity and pin count, the other signalling rate and power.

The procurement consequence is a sentence rather than an analysis, and it is the most immediately usable finding in this chapter for anyone who buys machines: \emph{specify terabytes per second per stack in the contract, not ``HBM4.''} An agreement written to the generation name can be filled, legally, at two-thirds of the bandwidth the buyer modeled.

What did \emph{not} improve is the strategically interesting half. The prefetch has been 256 bits since HBM1 in 2014, and the die density is 24\,Gb in both HBM3E and HBM4. HBM4's bandwidth gain---$1.6\times$ to $2.4\times$ over HBM3E depending on which vendor's bin is bought---therefore comes almost entirely from doubling the interposer I/O to 2,048 pins, with the doubling of banks per die to 256 reducing conflicts rather than adding raw bandwidth, and its capacity growth comes only from stack height~\cite{hc2026-hbm4-spread}. \emph{That is a packaging problem, not a DRAM problem.} It relocates the frontier off the memory cell---the layer on which CXMT's ${\sim}16$\,nm-class process is measurably behind and measurably closing---and onto the interposer, the bump pitch, the stack and the bonding step, which is the layer Section~\ref{sec:packaging} treats separately and Eq.~\eqref{eq:pkgyield} prices. For a country whose lithography is denied and whose packaging investment runs at RMB tens of billions a year, that relocation is on its face favourable. It is also the layer at which the second change bites.

\subsection{The base die became a logic die}

From HBM4 the base die of an HBM stack is fabricated on a 4\,nm foundry logic process [P]~\cite{hc2026-chbm}. That single fact is the inversion. A memory vendor buying leading-edge foundry logic capacity to build the bottom die of its own product is a chip designer, and Samsung's roadmap states what it intends to design. Phase one reclaims accelerator area by absorbing the memory controller into the base die. Phase two fills the freed base-die area with repair, telemetry, test and processing elements. Phase three deletes the interposer entirely, bonding the cell-array stack directly onto the accelerator---``zHBM,'' claimed at a 70\% reduction in I/O energy and 230\% of HBM4E bandwidth, with the ${\sim}100$\,W of DRAM power saved re-spent as an 8.3\% rise in accelerator power against a 1,200\,W baseline [R; vendor roadmap, unaudited]~\cite{hc2026-chbm}. And the commercial form is stated rather than inferred: \emph{custom HBM---``cHBM''---is explicitly not JEDEC-standardized}: a customer-specific logic base die on that 4\,nm process, sharing standard cell-array stacks. The two names are not interchangeable, and this book keeps them apart: cHBM is the phase-one and phase-two product, the co-designed logic base die; zHBM is the phase-three concept that deletes the interposer altogether by bonding the cell-array stack directly onto the xPU. Strategically, the memory vendor comes to own an interface that belonged to the accelerator designer.

The consequence this book must draw is not the one the deck draws. Samsung's framing is technical---area, energy, bandwidth. The structural change is that an HBM generation stops being a parameter set any qualified buyer can order against and becomes a per-customer product negotiated between a memory vendor and an accelerator team. That is good for large buyers with silicon teams---the hyperscalers, the two merchant GPU vendors, the handful of ASIC programs with a Broadcom or a Marvell attached---and bad for everyone else, national laboratories buying at merchant volumes very much included. The point is worth stating plainly rather than diplomatically: a centre procuring a few thousand accelerators does not get a logic die designed for it, and will increasingly buy the residue of a market whose leading products were co-designed for somebody else.

Score that against the ledger's ``three to four years behind,'' because a moving generational target is one thing and a market that has stopped having generations is another. Two readings are live and the public evidence does not separate them. The pessimistic one for China: catching HBM3E in 2027 no longer catches the frontier, because the frontier part is not a generation but a relationship---a logic die co-designed with an accelerator team---and access to accelerator design teams is precisely what the controls deny. On that reading the gap stops being measured in years of DRAM process, where it was closing, and starts being measured in bilateral commercial relationships China is not party to, where it cannot close at all. The optimistic one: a JEDEC generation is a target and a set of private bilateral products is not, so the \emph{merchant} HBM part---the one anyone can buy---stops advancing at the frontier rate once the fastest bins and the custom base dies are absorbed into co-designed products first. If the merchant market slows, the target CXMT is chasing slows with it. And China already runs the custom-HBM pattern inside a single national vertical: Huawei's HiBL~1.0 and HiZQ~2.0 are an accelerator designer specifying its own memory to an undisclosed domestic fab partner, which is structurally what Samsung now proposes to sell to everyone else. Vertical integration, the thing sanctions forced on the domestic ecosystem, is on this axis an advantage.

The chapter's finding is therefore restated rather than withdrawn, and the restatement is longer because it must now name two clocks: \emph{Chinese HBM is three to four years behind the merchant part, and outside the co-design relationships that increasingly define the frontier part}---the first gap closing at a measurable rate, the second not measurable at all from public evidence, and neither settled by the other [A].

\section{Packaging: the chokepoint hiding inside the memory ledger}
\label{sec:packaging}

Advanced packaging deserves separation from the HBM discussion, because it is a distinct chokepoint with distinct suppliers and a distinct clock. Producing logic and producing memory are insufficient if a country cannot \emph{assemble} them at yield and volume: through-silicon vias, hybrid bonding, interposers, large organic substrates, CoWoS-class capacity, known-good-die test, thermal interfaces, and increasingly optical and SerDes packaging are all separate industrial competences, and OSAT scaling is its own capital cycle. China's position here is markedly stronger than in lithography and markedly weaker than in assembly labour: Innotron's \$2.4B TSV facility, Tongfu's RMB 4.4B raise, more than RMB 27B of packaging expansion in H1 2026, and a domestic HBM assembly toolchain---TSV etch, hybrid bonding, precision placement---targeted at end-2026 HBM3~\cite{packaging,hbm-toolchain}. The strategic reading is that packaging is the layer where the domestic ecosystem is closing fastest \emph{and} the layer whose failure would strand both the logic and the memory investments; it is therefore the highest-variance item in Table~\ref{tab:semi-scoreboard}.

The 2026 disclosure record sharpens that judgement in three specific places, all of them from SK~hynix's own packaging material and all physical rather than commercial [P]~\cite{hc2026,hc2026-hbm4-spread}. First, stack height has become a warpage problem: going from twelve dies to sixteen takes package height only from 720 to 775\,\textmu m, so the die thickness falls to about $0.9\times$ and the inter-die gap roughly halves; beyond twenty dies the thermocompression flow does not survive at all and hybrid bonding becomes mandatory rather than optional. Second, the PHY power density rises from 0.5\,W/mm$^2$ at what SK~hynix labels sHBM4E to over 2.0\,W/mm$^2$ at sHBM5, crossing a line the vendor's own slide marks ``system limit'' at exactly the HBM4 generation---the interface, not the array, is what runs out of thermal budget. Third, and most consequential for yield economics, the assembly order has inverted: the memory stack is now attached before the rest of the package is complete, so an HBM defect scraps the entire system-in-package rather than a stack---and no vendor published a yield figure of any kind. Eq.~\eqref{eq:pkgyield} supplies the arithmetic; what changed is that the terms now multiply in an order which makes the memory yield the most expensive one to get wrong.

For the domestic program the reading is mixed and the two halves point opposite ways. The favourable half is that the frontier has moved onto the layer where China is investing hardest and where the tooling is least dependent on the lithography it cannot buy. The unfavourable half is sharper: the least mature item on the domestic HBM assembly list---Maxwell's hybrid bonding, targeted at end-2026 HBM3---is exactly the tool that stack heights beyond sixteen make compulsory, so the domestic schedule's most speculative element is the one the next generation needs first~\cite{hbm-toolchain}. Packaging remains the highest-variance row in Table~\ref{tab:semi-scoreboard}, and on this evidence the variance has widened rather than narrowed.

\section{A sanctions-elasticity model}
\label{sec:elasticity}

This book has used ``chokepoint decay'' as a descriptive phrase. It can be made semi-quantitative, and doing so is the honest way to compare controls that behave very differently. For each controlled input, write $\Delta t_{\text{eff}}$ for the \emph{effective delay}: the number of years by which a control postpones the adversary's attainment of the capability it targets, relative to a no-control counterfactual. It is not a formula but a function of five terms, each of which subtracts from the delay the control's designers assumed they were buying:
\begin{align*}
\Delta t_{\text{eff}} \;=\; \Phi\big(&\underbrace{\text{stockpile depth}}_{\text{years of inventory held}},\ \underbrace{\text{diversion}}_{\text{share obtainable anyway}},\ \underbrace{\text{domestic substitution rate}}_{\text{years to a home-grown equivalent}},\\
 &\underbrace{\text{architectural redesign}}_{\text{can the need be designed away?}},\ \underbrace{\text{yield penalty}}_{\text{cost of the inferior path}}\big),
\end{align*}
where $\Phi$ is left unspecified because no public dataset supports fitting it; what the decomposition buys is not a number but a checklist, and the terms are ordered so that the first three shorten the delay and the last two lengthen it. The terms have very different magnitudes by input. EUV lithography has no stockpile, negligible diversion at machine scale, a substitution rate measured in decades, and no architectural detour that fully substitutes---so the control buys real years. HBM has a large stockpile (exhausted by mid-2026), meaningful diversion, a substitution rate of three-plus years, and---critically---an \emph{architectural detour} in the capacity-first design of Chapter~\ref{ch:compute}, so the control buys less than its designers assumed. Advanced accelerators have shallow stockpiles, documented diversion at billion-dollar scale, a fast-improving domestic substitute, and a full architectural detour; the control has largely converted into a domestic market-reservation policy. EDA proved to have near-zero delay value: the ban was imposed and rescinded within six weeks, and domestic share doubled in the interval. Deposition, etch, and metrology sit in between.

The 2026 memory market exposes a sixth term the five omit, and it is not a technical one: the \emph{price of the substitute}. An architectural detour is only as cheap as the components it detours onto, and Section~\ref{sec:memory-wall} shows those components repricing by multiples inside eighteen months for reasons---HBM's die-area intensity, a decade of flat wafer starts, a two-year installation lag---that neither government controls and neither government's controls can reach. The corollary is uncomfortable but follows directly: a control's effective delay is partly a function of a global commodity price, so any $\Delta t_{\text{eff}}$ quoted without a date attached to it should be distrusted, and the HBM row in particular has a wider band in 2026 than the same analysis would have carried in 2024. Ranking controls by $\Delta t_{\text{eff}}$ rather than by strategic-sounding importance is, on this book's evidence, the single most useful analytical discipline available to export-control policy---and it would have predicted, in 2022, most of what has since happened.

\section{Equipment, EDA, materials: the long ladder}

The deepest dependency is the tool base, and here the trend, not the level, is the message. Naura reached \#5 in global semiconductor-equipment revenue in 2025 (from \#8 in 2022, revenue ${\sim}\$5.5$B, +31\%); AMEC (\#13) ships etch tools aimed at $\leq$5\,nm logic; three Chinese firms now sit in the global top twenty; domestic tools supply an estimated 20--30\% of Chinese fab demand, up from ${\sim}10\%$ three years prior~\cite{toolmakers-top20,naura}. The Huawei-linked, state-owned SiCarrier exhibited some thirty tool types within two years of founding and holds patents claiming 5\,nm-class SAQP-DUV flows~\cite{sicarrier}. Lithography remains the hard floor: SMEE ships at 90\,nm, 28\,nm-class immersion DUV is in trials (Yuliangsheng), and domestic EUV---LDP-source prototypes reportedly under test at Dongguan---is credibly dated no earlier than ${\sim}2030$ by sober assessments~\cite{litho}. The pattern of every denied layer repeats in EDA: the May 2025 US ban on EDA sales was rescinded within six weeks (traded against rare-earth flows---Chinese leverage from the previous wars, spent to protect the next one), but not before Empyrean shipped China's first full-process memory-EDA platform and domestic EDA share doubled~\cite{eda}. Silicon wafers: 3\% of 300\,mm capacity in 2020, ${\sim}28$--32\% in 2026~\cite{wafers}.

The money behind the ladder: Big Fund III alone is RMB 344B (${\sim}\$47.5$B) with a horizon to 2039, atop Funds I--II and provincial vehicles; China has been the world's largest equipment buyer (${\sim}40\%$ of global WFE) every year since 2020 and will remain so through 2027~\cite{bigfund,semi-wfe}. Self-sufficiency targets are missed and re-set---the ``Made in China 2025'' 70\% goal landed nearer 25--50\% depending on definition~\cite{self-sufficiency}---but the direction never reverses, and each US control action (December 2024's 140-entity action, the whiplash of the AI Diffusion Rule imposed January 2025 and rescinded May 2025) functions as demand guarantee for the domestic substitute~\cite{bis-dec24,policy-whiplash}. Meanwhile enforcement leaks at industrial scale: a \$2.5B smuggling scheme via Supermicro employees was charged in March 2026~\cite{smuggling}.

\section{The national-open logic at the silicon layer}

The semiconductor element lacks the visible ``openness'' branding of models and ISAs---one cannot open-source a fab. But the same national-open structure operates: open architecture above (RISC-V cores taped out on SMIC; XiangShan as shared national IP), unified standards across competing fabs and OSATs, state capital socializing the capex, and the deliberate choice of \emph{commodity} process technology---DUV multipatterning, LPDDR6, chiplet packaging---over the proprietary-premium branch (EUV, HBM4) wherever the commodity branch can be scaled domestically---a choice that Section~\ref{sec:hbm4-spread} shows has partly been made for it, because the premium branch has stopped being purchasable at all except by co-design. Table~\ref{tab:semi-scoreboard} summarizes the position.

\begin{table}[htbp]
  \centering
  \small
  \caption{The semiconductor element, mid-2026: chokepoints and their decay state.}
  \label{tab:semi-scoreboard}
  \setlength{\tabcolsep}{4pt}
  \begin{tabular}{p{39mm}p{48mm}p{51mm}}
    \toprule
    Chokepoint & Western position & Decay state \\
    \midrule
    Leading-edge logic & TSMC 2\,nm; EUV denied & SMIC 5\,nm-class DUV in volume; 60--80k wpm ramp; yield-limited, not capability-limited \\
    \addlinespace
    HBM, merchant part & SK hynix/Micron/Samsung lead; export-controlled; $1.5\times$ per-stack bandwidth spread inside one HBM4 generation & Potentially reduced as a sole dependency by the modeled capacity-first architecture [M]---not yet bypassed in any measured training system; domestic HBM2E + now-exhausted stockpiles bridge \\
    \addlinespace
    HBM, co-designed part & Custom logic base die on 4\,nm, explicitly not JEDEC-standardized; sold to buyers with silicon teams & No decay path visible: the barrier is a commercial relationship, not a process step [A]. Huawei's HiBL/HiZQ is the same pattern run inside one national vertical \\
    \addlinespace
    Advanced packaging & CoWoS-class capacity, TSV, hybrid bonding, known-good-die test; the layer HBM4's gains now come from & Closing fastest of any layer (Innotron, Tongfu, $>$RMB 27B in H1 2026, domestic TSV and hybrid-bonding toolchain at end-2026 HBM3) [P/A]---and the highest-variance row here, since $\geq$20-high stacks make the least mature domestic tool compulsory \\
    \addlinespace
    LPDDR6 & JEDEC-standard commodity class & CXMT roadmap targets a schedule comparable with leading suppliers [A/R]; qualified mass production not yet independently verified \\
    \addlinespace
    DRAM commodity & Samsung/SK/Micron oligopoly & CXMT \#4 and rising; swing supplier in shortage; pricing above Samsung \\
    \addlinespace
    Lithography & ASML monopoly & Hard floor; 28\,nm DUV trials; EUV ${\sim}2030$; mitigated by multipatterning + capacity \\
    \addlinespace
    Equipment (non-litho) & AMAT/Lam/TEL lead & Naura \#5 globally; 20--30\% domestic share and rising \\
    \addlinespace
    EDA & Synopsys/Cadence duopoly & Ban imposed and rescinded in 6 weeks; domestic share doubling from low base \\
    \bottomrule
  \end{tabular}
\end{table}

The element-level conclusion, stated at the evidence grades the table carries rather than above them: no single Chinese node matches its Western counterpart at the frontier, and the architecture of Chapters~\ref{ch:model} and~\ref{ch:compute} was chosen so that none needs to. But ``sufficient'' is a claim about an integrated system, and the honest formulation is conditional: \emph{the emerging domestic stack contains plausible paths to sufficiency, each supported to a different evidence level---logic capability [V] but logic economics [A], LPDDR roadmap [A/R], packaging capability [P/A], software maturity [M]---and an end-to-end frontier training run remains the decisive integration test that none of the paths has yet passed}. Logic yield, packaging economics, memory qualification, software maturity, and LPDDR training performance have never been demonstrated \emph{simultaneously in one run}; the claim register carries that run as CR-09, the ``Mate~60 moment'' of AI. What the table does establish unconditionally is the strategic geometry: the chokepoint strategy assumed China would try to rebuild the Western machine and could be stopped at its scarcest parts. China is building a different machine---whether the different machine works end to end is the open question this book's forecasts price at 65\% within 2--4 years. What 2026 added to that question is a variable neither side chose and neither side controls: the price of memory, which moved by multiples in eighteen months, which prices both machines, and which is now the largest single source of uncertainty in every cost comparison this book makes.
